\scanpage{36}

% B01 | printed letterhead, upper left
% letterhead
\begin{center}
UNIVERSITÉ DE PARIS\\
---$\cdot$---\\
ÉCOLE NORMALE SUPÉRIEURE\\
---$\cdot$---\\
45, RUE D'ULM\\
ODÉON 06$\cdot$45\\
---$\cdot$---
\end{center}

% B02 | first three lines, crossed out as a block
\deleted{Soit $E$ \deleted{\unsure{un ensemble simplicial}} un C-w-complexe, $G$ un groupe discret, opérant sur $E$ par des transformations qui préservent la décomposition cellulaire, $B=E/G$ \deleted{l'ensemble} le C-w-complexe quotient de $E$}

% B03 | Lemme 1, two lines
\uline{Lemme 1} : $\forall x$ sommet de $E$, $\beta$ chemin de $B$ tel que $\partial_1(\beta)=p(x)$, $\exists\gamma$ chemin de $E$ tel que $\partial_1(\gamma)=x$ et $p(\gamma)=\beta$

% B04 | Dém, three lines
\uuline{Dém} Il suffit de le voir si $\beta$ est une arête. Soit $\lambda$ une arête de $B$ telle que $p\lambda=\beta$. On a \deleted{$\unsure{\partial}$} $p(\partial_1\lambda)=\partial_1(p\lambda)=\partial_1\beta=p(x)$, d'où $\exists g\deleted{\unsure{g}}\in G$ $x=g.\partial_1\lambda$, et $\gamma=g\lambda$ répond à la question

% B05 | Lemme 2, below centre fold
\uline{Lemme 2} : l'image de $p_*:\pi_1(E)\to\pi_1(B)$ est un sous groupe distingué de $\pi_1(B)$

% B06 | Dém, four lines
\uuline{Dém} : Soit $\gamma$ un lacet de \deleted{$\unsure{B}$} $E$, $\beta$ un lacet de $B$, $\lambda$ un chemin de $E$ tel que $\partial_1(\lambda)=x_0$ et $p_{\deleted{\unsure{*}}}(\lambda)=\beta$. On a $p(\partial_0\lambda)=\partial_0\beta=b_0=px_0$, d'où $\partial_0\lambda=gx_0$, et le lacet $(\lambda,g\gamma,-\lambda)$ se projette en $(\beta,p(\gamma),-\beta)$ ce qui démontre le lemme 2.

% B07 | underlined heading in lower quarter
\uline{Définition de $\chi:G\longrightarrow\pi_1(B)/\operatorname{Im}\pi_1(E)$ quand $E$ connexe.}

% B08 | final paragraph, five lines
Pour tout $g\in G$, soit $\gamma$ un chemin de $E$ tel que $\partial_1(\gamma)=x_0$, $\partial_0(\gamma)=gx_0$. \deleted{\unsure{soit $\chi(g)$ la classe de|et $\chi(g)$ la classe de}} $p(\gamma)$ est un lacet de $B$, soit $\chi(g)$ sa classe dans $\pi_1(B)/\operatorname{im}\pi_1(E)$. Elle ne dépend pas du chemin $\gamma$ choisi, car, si $\gamma'$ est un autre chemin, $(\gamma,-\gamma')$ est un lacet de $B$, et \deleted{\unsure{$p(\gamma)$ est homotope|par définition}} $p(\gamma)$ et $p(\gamma')$ ont même classe mod l'image de $\pi_1(E)$.

% archive sheet number: 35

